\section{Final thoughts: does Shor’s algorithm really break RSA?}
During your lecture, you went through all the Shor’s algorithm, you had a whole journey with it, and I hope that I was able to make you enjoy the journey. You probably have a lot of questions in your mind now, and one of them would be that you saw with your own eyes that we have an algorithm that breaks our usual encryption. So how is our data safe?\nl
Actually, we not really did. We are assuming that we have a quantum computer that is able to run Shor’s algorithm, and we don’t. Shor’s algorithm would need 4100 logical qubits, or 20 000 000 physical qubits to run.\nl
Currently, the main actor in this field is IBM, and they are aiming for 1121 physical qubits soon, far from our 20 000 000 ones. But the day we will be able to make a quantum computer with 20 000 000 physical qubits, we will definitely be able to break RSA, so the question is not if we are going to be able to, but \textbf{when}? 
